%%%%%%%%%%  scan idx 345  |  folio 332  %%%%%%%%%%
\textbf{Proposition 3.11.} Soient \(A\) un \(\lambda\)-anneau, et \(y_i\), \(i=1,\dots,n\) des éléments de \(A\) tels que \(\lambda^j(y_i)=0\) pour \(j\ge2\). On pose \(x=\sum_i y_i\)~; alors :
\begin{enumerate}[label=\roman*)]
\item pour \(k>n\), on a \(\gamma^k(x-n)=0\)
\item pour \(1\le k\le n\), \(\gamma^k(x-n)\) est la \(k\)-ième fonction symétrique des \(y_i-1\).
\end{enumerate}
Soit \(j\ge2\). On a alors \(\gamma^j(y_i-1)=\lambda^j(y_i+j-2)=\sum_{p+q=j}\lambda^p(y_i)\lambda^q(j-2)\). Comme \(\lambda^p(y_i)=0\) pour \(p\ge2\), et \(\lambda^q(j-2)=0\) pour \(q>j-2\), \(\gamma^j(y_i-1)=0\). La proposition résulte alors de ce que \(\gamma_t(x-n)=\prod_{i=1}^{n}\gamma_t(y_i-1)\).

\section*{4. \(\lambda\)-anneaux engendrés par générateurs et relations}

On se place toujours sur \(\mathcal C_0\).

\textbf{4.1.} Soit \(K\) un \(\lambda\)-anneau. On considère une famille d'indéterminées \((X_i)_{i\in I}\), et soit \(A=K[X_i,\lambda^n X_i]_{i\in I,\ n\ge2}\)~; on veut définir sur \(A\) une structure de \(\lambda\)-anneau prolongeant celle de \(K\).

Pour se donner un homomorphisme \(\lambda_t:A\to\Ghat(A)\), il suffit de se donner les images des \(X_i\), \(\lambda^n X_i\), ce qui définit un homomorphisme d'anneaux : \(A\to\Ghat(A)_o\), et pour qu'il définisse une structure de \(\lambda\)-anneau sur \(A\), il faut vérifier la commutativité du diagramme (3.7.1). On pose donc :
\begin{equation}\tag{4.1.1}
  \lambda_t(X_i)=1+X_it+\sum_{n\ge2}\lambda^n X_i\,t^n~;
  \qquad
  \lambda_t(\lambda^n X_i)=\lambda^n(\lambda_t(X_i)).
\end{equation}
Le diagramme (3.7.1) étant un diagramme d'homomorphismes d'anneaux, il suffit de vérifier sa commutativité pour les \(X_i\), \(\lambda^n X_i\) (elle est connue par hypothèse sur les éléments de \(K\)), ce qui résulte facilement des définitions (4.1.1) et de ce que \(\lambda_u\) et \(\Ghat^u(\lambda_t)\) sont des \(\lambda\)-homomorphismes.

